\documentclass[12pt]{article}

\usepackage[T1]{fontenc}
\usepackage[utf8]{inputenc}
\usepackage[brazil]{babel}
\usepackage{lmodern}
\usepackage{microtype}
\usepackage[a4paper,margin=2.5cm]{geometry}
\usepackage{setspace}
\usepackage{indentfirst}

\setstretch{1.15}

\title{O Macroeconomista como Cientista e Engenheiro \thanks{Texto Traduzido livremente por Luiz Mario Andrade com o auxílio da ferramenta Chat GPT 5.4}}
\author{N. Gregory Mankiw}
\date{Outono de 2006}

\begin{document}

\maketitle

Os economistas gostam de assumir a postura de cientistas. Eu sei disso porque frequentemente o faço. Quando ensino alunos de graduação, descrevo conscientemente o campo da economia como uma ciência, para que nenhum estudante comece o curso pensando que está embarcando em algum empreendimento acadêmico vago. Nossos colegas no departamento de física do outro lado do campus podem achar divertido que os vejamos como primos próximos, mas somos rápidos em lembrar a quem quiser ouvir que os economistas formulam teorias com precisão matemática, coletam enormes conjuntos de dados sobre o comportamento individual e agregado e exploram as técnicas estatísticas mais sofisticadas para chegar a julgamentos empíricos isentos de viés e ideologia (ou assim gostamos de pensar).

Tendo passado recentemente dois anos em Washington como conselheiro econômico em um momento em que a economia dos EUA lutava para sair de uma recessão, lembro-me de que o subcampo da macroeconomia nasceu não como uma ciência, mas mais como um tipo de engenharia. Deus colocou os macroeconomistas na Terra não para propor e testar teorias elegantes, mas para resolver problemas práticos. Os problemas que Ele nos deu, além disso, não foram de dimensões modestas. O problema que deu origem ao nosso campo — a Grande Depressão da década de 1930 — foi uma retração econômica de escala sem precedentes, incluindo rendas tão deprimidas e desemprego tão disseminado que não é exagero dizer que a viabilidade do sistema capitalista foi posta em questão.

Este ensaio oferece uma breve história da macroeconomia, juntamente com uma avaliação do que aprendemos. Minha premissa é que o campo evoluiu através dos esforços de dois tipos de macroeconomistas — aqueles que entendem o campo como um tipo de engenharia e aqueles que gostariam que ele fosse mais uma ciência. Engenheiros são, antes de tudo, solucionadores de problemas. Em contraste, o objetivo dos cientistas é entender como o mundo funciona. A ênfase da pesquisa dos macroeconomistas variou ao longo do tempo entre esses dois motivos. Enquanto os primeiros macroeconomistas eram engenheiros tentando resolver problemas práticos, os macroeconomistas das últimas décadas estiveram mais interessados em desenvolver ferramentas analíticas e estabelecer princípios teóricos. Essas ferramentas e princípios, no entanto, têm demorado a encontrar seu caminho nas aplicações. À medida que o campo da macroeconomia evoluiu, um tema recorrente é a interação — às vezes produtiva e às vezes não — entre os cientistas e os engenheiros. A desconexão substancial entre a ciência e a engenharia da macroeconomia deve ser um fato humilhante para todos nós que trabalhamos na área.

Para evitar qualquer confusão, devo dizer logo de início que a história que conto não é uma história de mocinhos e vilões. Nem cientistas nem engenheiros têm uma reivindicação de maior virtude. A história também não é sobre pensadores profundos e encanadores simplórios. Professores de ciências geralmente não são melhores em resolver problemas de engenharia do que professores de engenharia em resolver problemas científicos. Em ambos os campos, os problemas de ponta são problemas difíceis, bem como intelectualmente desafiadores. Assim como o mundo precisa de cientistas e engenheiros, ele precisa de macroeconomistas de ambas as mentalidades. Mas acredito que a disciplina avançaria de forma mais suave e proveitosa se os macroeconomistas sempre tivessem em mente que seu campo tem um papel duplo.

\section{A Revolução Keynesiana}

A palavra "macroeconomia" aparece pela primeira vez na literatura acadêmica na década de 1940. Certamente, os tópicos da macroeconomia — inflação, desemprego, crescimento econômico, ciclo de negócios e política monetária e fiscal — há muito intrigam os economistas. No século XVIII, por exemplo, David Hume (1752) escreveu sobre os efeitos de curto e longo prazo das injeções monetárias; em muitos pontos, sua análise parece notavelmente semelhante ao que se poderia ver de um economista monetário moderno ou de um banqueiro central. Em 1927, Arthur Pigou publicou um livro intitulado \textit{Industrial Fluctuations} que tentava explicar o ciclo econômico. No entanto, o campo da macroeconomia como uma área de investigação distinta e ativa surgiu à sombra da Grande Depressão. Não há nada como uma crise para focar a mente.

A Grande Depressão teve um impacto profundo sobre aqueles que a viveram. Em 1933, a taxa de desemprego dos EUA atingiu 25 por cento, e o PIB real estava 31 por cento abaixo de seu nível de 1929. Todas as flutuações subsequentes na economia dos EUA foram marolas em um mar calmo em comparação com esse tsunami. Ensaios autobiográficos de economistas proeminentes dessa era, como Lawrence Klein, Franco Modigliani, Paul Samuelson, Robert Solow e James Tobin, confirmam que a Depressão foi um evento motivador fundamental em suas carreiras (Breit e Hirsch, 2004).

A \textit{Teoria Geral} de John Maynard Keynes foi o ponto focal nas discussões profissionais sobre como entender esses desenvolvimentos. Todos os cinco ganhadores do Nobel confirmam isso por experiência própria. Tobin relata a seguinte reação em Harvard, onde era estudante no final da década de 1930 e início da de 1940: "O corpo docente sênior era em sua maioria hostil... O corpo docente mais jovem e os estudantes de pós-graduação estavam entusiasmados com o livro de Keynes". Como costuma acontecer, os jovens tiveram maior percepção do que os velhos sobre o impacto das novas ideias. Keynes empatou com Alfred Marshall como o economista mais citado em periódicos econômicos na década de 1930 e foi o segundo mais citado na década de 1940, depois de Hicks (Quandt, 1976). Essa influência persistiu por muitos anos. Keynes ocupou o 14º lugar em citações para o período de 1966 a 1986, embora tenha morrido duas décadas antes do início do período (Garfield, 1990).

A revolução keynesiana influenciou não apenas a pesquisa econômica, mas também a pedagogia. O livro didático clássico de Samuelson, \textit{Economics}, foi publicado pela primeira vez em 1948, e sua organização refletia sua percepção do que a profissão tinha a oferecer ao leitor leigo. Oferta e demanda, que hoje estão no centro de como ensinamos economia aos calouros, não foram introduzidas até a página 447 do livro de 608 páginas. A macroeconomia vinha primeiro, incluindo conceitos como o multiplicador da política fiscal e o paradoxo da parcimônia. Samuelson escreveu (na p. 253): "Embora grande parte desta análise se deva a um economista inglês, John Maynard Keynes... hoje seus fundamentos amplos são cada vez mais aceitos por economistas de todas as escolas de pensamento".

Quando um economista moderno lê a \textit{Teoria Geral}, a experiência é ao mesmo tempo estimulante e frustrante. Por um lado, o livro é o trabalho de uma mente brilhante sendo aplicada a um problema social cuja atualidade e enormidade não podem ser questionadas. Por outro lado, embora o livro seja extenso em sua análise, ele parece incompleto como questão de lógica. Muitos fios ficam soltos. O leitor continua perguntando: qual é, precisamente, o modelo econômico que une todas as peças?

Logo após Keynes publicar a \textit{Teoria Geral}, uma geração de macroeconomistas trabalhou para responder a essa pergunta, transformando sua grande visão em um modelo mais simples e concreto. Uma das primeiras e mais influentes tentativas foi o modelo IS-LM proposto por John Hicks (1937), então com 33 anos. Franco Modigliani (1944), com 26 anos, estendeu e explicou o modelo mais plenamente. Até hoje, o modelo IS-LM continua sendo a interpretação de Keynes oferecida nos livros didáticos de macroeconomia de nível intermediário mais amplamente utilizados. Alguns críticos keynesianos do modelo IS-LM queixam-se de que ele simplifica demais a visão econômica oferecida por Keynes na \textit{Teoria Geral}. Até certo ponto, isso pode ser verdade. O objetivo do modelo era simplificar uma linha de argumentação que, de outra forma, seria difícil de seguir. A linha entre simplificar e simplificar demais muitas vezes não é clara.

Enquanto teóricos como Hicks e Modigliani desenvolviam modelos keynesianos adequados para o quadro negro da sala de aula, econometricistas como Klein trabalhavam em modelos mais aplicados que podiam ser levados aos dados e usados para análise de políticas. Com o tempo, na esperança de se tornarem mais realistas, os modelos tornaram-se maiores e acabaram incluindo centenas de variáveis e equações. Na década de 1960, havia muitos modelos concorrentes, cada um baseado na contribuição de keynesianos proeminentes da época, como o modelo Wharton associado a Klein, o modelo DRI (Data Resources, Inc.) associado a Otto Eckstein e o modelo MPS (MIT–Penn–Social Science Research Council) associado a Albert Ando e Modigliani. Esses modelos foram amplamente utilizados para previsão e análise de políticas. O modelo MPS foi mantido pelo Federal Reserve por muitos anos e se tornaria o precursor do modelo FRB/US, que ainda é mantido e usado pela equipe do Fed.

Embora esses modelos diferissem em detalhes, suas semelhanças eram mais marcantes do que suas diferenças. Todos tinham uma estrutura essencialmente keynesiana. No fundo da mente de cada construtor de modelos estava o mesmo modelo ensinado aos estudantes de graduação hoje: uma curva IS relacionando as condições financeiras e a política fiscal aos componentes do PIB, uma curva LM que determinava as taxas de juros como o preço que equilibra a oferta e a demanda de moeda, e algum tipo de curva de Phillips que descreve como o nível de preços responde ao longo do tempo às mudanças na economia.

Como questão de ciência, a \textit{Teoria Geral} foi um sucesso notável. A revolução que inspirou atraiu muitas das melhores mentes jovens de seu tempo. Sua produção prodigiosa ofereceu uma nova maneira de entender as flutuações econômicas de curto prazo. Refletindo sobre esses eventos, Samuelson (1988) ofereceu um resumo sucinto: "A revolução keynesiana foi o evento mais significativo na ciência econômica do século XX". Esse sentimento é compartilhado por muitos economistas de sua geração.

No entanto, a revolução keynesiana não pode ser entendida meramente como um avanço científico. Em grande medida, Keynes e os construtores de modelos keynesianos tinham a perspectiva de engenheiros. Eles eram motivados por problemas no mundo real e, uma vez desenvolvidas suas teorias, estavam ansiosos para colocá-las em prática. Até sua morte em 1946, o próprio Keynes esteve fortemente envolvido em oferecer conselhos de política. O mesmo ocorreu com os primeiros keynesianos americanos. Tobin, Solow e Eckstein tiraram tempo de suas atividades acadêmicas durante a década de 1960 para trabalhar no Conselho de Consultores Econômicos. O corte de impostos de Kennedy, finalmente aprovado em 1964, foi em muitos aspectos o resultado direto do consenso keynesiano emergente e dos modelos que o encarnavam.

\section{Os Novos Clássicos}

No final da década de 1960, rachaduras no consenso keynesiano estavam começando a aparecer. Essas rachaduras cresceriam e se tornariam fissuras que acabariam por desmoronar o consenso macroeconômico e minar a confiança nos modelos econométricos convencionais. Em seu lugar, uma visão mais clássica da economia ressurgiria.

A primeira onda da nova economia clássica foi o monetarismo, e seu proponente mais notável foi Milton Friedman. O trabalho inicial de Friedman (1957) sobre a hipótese da renda permanente não era diretamente sobre moeda ou o ciclo econômico, mas certamente tinha implicações para a teoria do ciclo econômico. Foi, em parte, um ataque à função de consumo keynesiana, que fornecia a base para os multiplicadores de política fiscal que eram centrais para a teoria e as prescrições de política keynesianas. Se a propensão marginal a consumir a partir da renda transitória é pequena, como a teoria de Friedman sugeria, então a política fiscal teria um impacto muito menor na renda de equilíbrio do que muitos keynesianos acreditavam.

A \textit{Monetary History of the United States} de Friedman e Schwartz (1963) estava mais diretamente preocupada com o ciclo econômico e também minou o consenso keynesiano. A maioria dos keynesianos via a economia como inerentemente volátil, constantemente fustigada pelos "espíritos animais" mutáveis dos investidores. Friedman e Schwartz sugeriram que a instabilidade econômica deveria ser atribuída não a atores privados, mas sim a uma política monetária inepta. A implicação era que os formuladores de políticas deveriam se dar por satisfeitos se não causassem danos seguindo regras de política simples. Embora a regra proposta por Friedman de crescimento constante nos agregados monetários tenha poucos adeptos hoje, foi um precursor inicial dos regimes de metas de inflação agora em vigor em muitos dos bancos centrais do mundo.

O Discurso Presidencial de Friedman para a American Economic Association em 1968, juntamente com Phelps (1968), mirou o elo mais fraco no modelo keynesiano: o \textit{trade-off} da curva de Phillips entre inflação e desemprego. Pelo menos desde Samuelson e Solow (1960), algum tipo de curva de Phillips fazia parte do consenso keynesiano, mesmo que não fosse uma visão endossada pelo próprio Keynes. Samuelson e Solow (1960) entendiam a fragilidade teórica desse \textit{trade-off}, e seu artigo estava cheio de ressalvas sobre por que o \textit{trade-off} de curto e longo prazo poderia diferir. Mas a literatura subsequente esqueceu essas ressalvas com muita facilidade. A curva de Phillips fornecia uma maneira conveniente de completar o modelo keynesiano, que sempre tinha dificuldade em explicar por que os preços não equilibravam os mercados e como o nível de preços se ajustava ao longo do tempo.

Friedman argumentou que o \textit{trade-off} entre inflação e desemprego não se manteria no longo prazo, quando os princípios clássicos deveriam se aplicar e a moeda deveria ser neutra. O \textit{trade-off} aparecia nos dados porque, no curto prazo, a inflação é frequentemente imprevista e a inflação imprevista pode reduzir o desemprego. O mecanismo particular que Friedman sugeriu foi a ilusão monetária por parte dos trabalhadores. Mais importante para o desenvolvimento da macroeconomia foi o fato de Friedman colocar as expectativas no centro do palco.

Isso preparou o caminho para a segunda onda da nova economia clássica — a revolução das expectativas racionais. Em uma série de artigos altamente influentes, Robert Lucas estendeu o argumento de Friedman. Em seu "Econometric Policy Evaluation: A Critique", Lucas (1976) argumentou que os modelos keynesianos convencionais eram inúteis para a análise de políticas porque não levavam as expectativas a sério; como resultado, as relações empíricas estimadas que compunham esses modelos quebrariam se uma nova política fosse implementada. Lucas (1973) também propôs uma teoria do ciclo econômico baseada nas suposições de informação imperfeita, expectativas racionais e equilíbrio de mercado. Nessa teoria, a política monetária importa apenas na medida em que surpreende as pessoas e as confunde sobre os preços relativos. Barro (1977) ofereceu evidências de que esse modelo era consistente com os dados de séries temporais dos EUA. Sargent e Wallace (1975) apontaram uma implicação política fundamental: como é impossível surpreender pessoas racionais sistematicamente, a política monetária sistemática visando estabilizar a economia está fadada ao fracasso.

A terceira onda da nova economia clássica foram as teorias dos ciclos reais de negócios de Kydland e Prescott (1982) e Long e Plosser (1983). Assim como as teorias de Friedman e Lucas, estas foram construídas sobre a suposição de que os preços se ajustam instantaneamente para equilibrar os mercados — uma diferença radical em relação à teorização keynesiana. Mas, ao contrário de seus predecessores novos clássicos, as teorias dos ciclos reais de negócios omitiam qualquer papel da política monetária, antecipada ou não, na explicação das flutuações econômicas. A ênfase mudou para o papel de choques aleatórios na tecnologia e a substituição intertemporal no consumo e no lazer que esses choques induziam.

Como resultado das três ondas da nova economia clássica, o campo da macroeconomia tornou-se cada vez mais rigoroso e cada vez mais ligado às ferramentas da microeconomia. Os modelos de ciclos reais de negócios eram exemplos específicos e dinâmicos da teoria do equilíbrio geral de Arrow-Debreu. De fato, este era um de seus principais pontos de venda. Com o tempo, os defensores deste trabalho recuaram da suposição de que o ciclo econômico é impulsionado por forças reais em oposição a monetárias, e começaram a enfatizar as contribuições metodológicas deste trabalho. Hoje, muitos macroeconomistas vindos da tradição nova clássica aceitam alegremente a suposição keynesiana de preços rígidos, desde que essa suposição esteja inserida em um modelo adequadamente rigoroso no qual os atores econômicos são racionais e olham para o futuro. Devido a essa mudança de ênfase, a terminologia evoluiu, e essa classe de trabalho agora é frequentemente designada pelo rótulo de teoria do "equilíbrio geral estocástico dinâmico". Mas estou me adiantando na história.

Na época em que as três ondas novas clássicas estavam chegando à costa nas décadas de 1970 e 1980, um de seus objetivos era minar os antigos modelos macroeconométricos keynesianos tanto como uma questão de ciência quanto de engenharia. Em seu artigo "After Keynesian Macroeconomics", Lucas e Sargent (1979) escreveram: "Para a política, o fato central é que as recomendações de política keynesianas não têm base mais sólida, num sentido científico, do que as recomendações de economistas não keynesianos ou, aliás, de não economistas". Embora Sargent e Lucas achassem que a engenharia keynesiana se baseava em uma ciência defeituosa, eles sabiam que a escola nova clássica (por volta de 1979) ainda não tinha um modelo pronto para levar a Washington: "Consideramos os melhores modelos de equilíbrio atualmente existentes como protótipos de modelos futuros e melhores que, esperamos, serão de uso prático na formulação de políticas". Eles também arriscaram que tais modelos estariam disponíveis "em dez anos, se tivermos sorte". Voltarei mais tarde à questão de se essa perspectiva se concretizou como eles esperavam.

Como sugerem essas citações, aqueles engajados no movimento novo clássico não foram tímidos sobre suas intenções nem modestos sobre suas realizações. Lucas ofereceu uma avaliação ainda mais contundente em um artigo de 1980 intitulado "The Death of Keynesian Economics": "Não se consegue encontrar bons economistas com menos de quarenta anos que se identifiquem ou identifiquem seu trabalho como 'keynesiano'. De fato, as pessoas até se ofendem se forem referidas como 'keynesianas'. Em seminários de pesquisa, as pessoas não levam mais a sério a teorização keynesiana; o público começa a sussurrar e a rir uns com os outros". No entanto, no exato momento em que Lucas escrevia o elogio fúnebre para a economia keynesiana, a profissão estava prestes a acolher uma geração de "novos keynesianos".

\section{Os Novos Keynesianos}

Os economistas atraídos pela abordagem keynesiana do ciclo econômico há muito se sentiam desconfortáveis com a questão dos microfundamentos. De fato, um artigo de 1946 de Lawrence Klein, um dos primeiros a usar o termo "macroeconomia", começa assim: "Muitos dos modelos matemáticos recém-construídos de sistemas econômicos, especialmente as teorias do ciclo econômico, estão muito vagamente relacionados ao comportamento de famílias ou empresas individuais, que devem formar a base de todas as teorias do comportamento econômico". Todos os economistas modernos são, até certo ponto, clássicos. Todos ensinamos aos nossos alunos sobre otimização, equilíbrio e eficiência de mercado. Como reconciliar essas duas visões da economia — uma fundamentada na mão invisível de Adam Smith e nas curvas de oferta e demanda de Alfred Marshall, a outra fundamentada na análise de Keynes de uma economia sofrendo de demanda agregada insuficiente — tem sido uma questão profunda e persistente desde que a macroeconomia começou como um campo de estudo separado.

Os primeiros keynesianos, como Samuelson, Modigliani e Tobin, pensavam ter reconciliado essas visões no que às vezes é chamado de "síntese neoclássica-keynesiana". Esses economistas acreditavam que a teoria clássica de Smith e Marshall estava correta no longo prazo, mas a mão invisível poderia ficar paralisada no curto prazo descrito por Keynes. O horizonte temporal importava porque alguns preços — mais notavelmente o preço do trabalho — ajustavam-se lentamente ao longo do tempo. Os primeiros keynesianos acreditavam que os modelos clássicos descreviam o equilíbrio para o qual a economia evoluía gradualmente, mas que os modelos keynesianos ofereciam a melhor descrição da economia em qualquer momento no tempo em que os preços fossem razoavelmente considerados pré-determinados.

A síntese neoclássica-keynesiana é coerente, mas também é vaga e incompleta. Enquanto os novos economistas clássicos responderam a esses defeitos rejeitando a síntese e começando do zero, os novos economistas keynesianos pensaram que havia muito a preservar. Seu objetivo era usar as ferramentas da microeconomia para dar maior precisão ao compromisso desconfortável alcançado pelos primeiros keynesianos. A síntese neoclássica-keynesiana era como uma casa construída na década de 1940: os novos clássicos olharam para seus sistemas desatualizados e concluíram que era uma demolição total, enquanto os novos keynesianos admiraram o artesanato do velho mundo e o abraçaram como uma oportunidade para uma grande reforma.

A primeira onda de pesquisa que pode ser chamada de "nova keynesiana" é o trabalho sobre o desequilíbrio geral (Barro e Grossman, 1971; Malinvaud, 1977). Essas teorias visavam usar as ferramentas da análise de equilíbrio geral para entender a alocação de recursos que resulta quando os mercados não se equilibram. Salários e preços eram tomados como dados. O foco estava em como a falha de um mercado em se equilibrar influencia a oferta e a demanda em mercados relacionados. De acordo com essas teorias, a economia pode se encontrar em um de vários regimes, dependendo de quais mercados estão experimentando excesso de oferta e quais estão experimentando excesso de demanda. O regime mais interessante — no sentido de corresponder melhor ao que observamos durante as contrações econômicas — é o chamado regime "keynesiano", no qual tanto o mercado de bens quanto o mercado de trabalho exibem excesso de oferta. No regime keynesiano, o desemprego surge porque a demanda de trabalho é muito baixa para garantir o pleno emprego aos salários vigentes; a demanda de trabalho é baixa porque as empresas não podem vender tudo o que desejam aos preços vigentes; e a demanda pela produção das empresas é inadequada porque muitos clientes estão desempregados. Recessões e depressões resultam de um círculo vicioso de demanda insuficiente, e um estímulo à demanda pode ter efeitos multiplicadores.

A segunda onda da pesquisa nova keynesiana visava explorar como o conceito de expectativas racionais poderia ser usado em modelos sem a suposição de equilíbrio de mercado. Até certo ponto, esse trabalho estava respondendo à conclusão de Sargent e Wallace (1975) sobre a irrelevância da política monetária, mostrando como a política monetária sistemática poderia potencialmente estabilizar a economia, apesar das expectativas racionais (Fischer, 1977). Até certo ponto, foi motivado pelo desejo de encontrar um modelo empiricamente realista de dinâmica inflacionária (Taylor, 1980). O calcanhar de Aquiles desse trabalho foi o fato de assumir uma forma de contratos de trabalho que, embora talvez justificável em bases empíricas, era difícil de conciliar com os princípios microeconômicos.

Como grande parte da tradição keynesiana se baseava na premissa de que os salários e preços não equilibram os mercados, a terceira onda de pesquisa nova keynesiana visava explicar por que esse era o caso. Várias hipóteses foram exploradas: que as empresas enfrentam "custos de menu" quando optam por mudar seus preços; que as empresas pagam a seus trabalhadores "salários de eficiência" acima do nível de equilíbrio de mercado para aumentar a produtividade do trabalhador; e que os definidores de salários e preços se desviam da racionalidade perfeita. Mankiw (1985) e Akerlof e Yellen (1985) apontaram que, quando as empresas têm poder de mercado, há grandes diferenças entre os cálculos de custo-benefício privados e sociais em relação ao ajuste de preços, de modo que um equilíbrio de preços rígidos poderia ser privadamente racional (ou quase racional) enquanto socialmente muito dispendioso. Blanchard e Kiyotaki (1987) mostraram que parte dessa divergência entre incentivos privados e sociais resulta de uma externalidade de demanda agregada: quando uma empresa corta seus preços, ela aumenta os saldos monetários reais e, assim, a demanda pelos produtos de todas as empresas. Ball e Romer (1990) estabeleceram que há uma forte complementaridade entre rigidezes reais e nominais, de modo que qualquer motivo para evitar mudanças de preços relativos exacerbaria a lentidão dos preços nominais.

Retrospectivamente, essas várias contribuições novas keynesianas estavam mais relacionadas e eram mais complementares do que pareciam na época, mesmo para as pessoas que trabalhavam nelas. Por exemplo, é tentador ver o trabalho inicial sobre desequilíbrio geral como um beco sem saída — um programa de pesquisa que semeou as sementes de sua própria destruição por sua suposição de preços pré-determinados. De fato, esse trabalho raramente encontra seu caminho nas listas de leitura hoje. No entanto, também se pode ver uma progressão de ideias relacionadas sobre como a economia funciona quando os preços não se movem instantaneamente para equilibrar a oferta e a demanda.

Existe, por exemplo, uma relação interessante, mas raramente notada, entre a primeira e a terceira ondas da economia nova keynesiana. Em particular, pode-se ver a terceira onda como estabelecendo a centralidade do regime keynesiano destacado na primeira onda. Quando as empresas têm poder de mercado, elas cobram preços acima do custo marginal, por isso sempre querem vender mais aos preços vigentes. Em certo sentido, se todas as empresas têm algum grau de poder de mercado, então os mercados de bens estão tipicamente em um estado de excesso de oferta. Esta teoria do mercado de bens é muitas vezes casada com uma teoria do mercado de trabalho com salários acima do equilíbrio, como o modelo de salário de eficiência. Nesse caso, o regime "keynesiano" de excesso de oferta generalizado não é apenas um resultado possível para a economia, mas o típico.

Em meu julgamento, essas três ondas de pesquisa nova keynesiana somaram-se a uma teoria microeconômica coerente para a falha da mão invisível em funcionar para fenômenos macroeconômicos de curto prazo. Entendemos como os mercados interagem quando há rigidez de preços, o papel que as expectativas podem desempenhar e os incentivos que os definidores de preços enfrentam ao escolher se mudam ou não os preços. Como questão de ciência, houve muito sucesso nesta pesquisa (embora, como participante, eu não possa pretender ser inteiramente objetivo). O trabalho não foi revolucionário, mas não tentava ser. Em vez disso, foi contrarrevolucionário: seu objetivo era defender a essência da síntese neoclássica-keynesiana do ataque novo clássico.

Este trabalho também foi bem-sucedido como uma questão de engenharia? Ajudou os formuladores de políticas a conceber melhores políticas para lidar com o ciclo econômico? O julgamento aqui deve ser menos positivo — um tópico ao qual voltarei em breve.

Mas é notável que os novos keynesianos eram, por temperamento, mais inclinados a se tornarem engenheiros macroeconômicos do que os economistas que trabalhavam na tradição nova clássica. Entre os líderes da escola nova clássica, nenhum (até onde eu sei) deixou a academia para assumir um cargo significativo na política pública. Em contraste, o movimento novo keynesiano, como a geração anterior de keynesianos, estava cheio de pessoas que trocaram alguns anos na torre de marfim por uma estadia na capital do país. Exemplos incluem Stanley Fischer, Larry Summers, Joseph Stiglitz, Janet Yellen, John Taylor, Richard Clarida, Ben Bernanke e eu mesmo. Os quatro primeiros desses economistas foram para Washington durante os anos Clinton; os quatro últimos durante os anos Bush. A divisão dos economistas entre novos clássicos e novos keynesianos não é, fundamentalmente, entre a direita política e a esquerda política. Em maior medida, é uma divisão entre cientistas puros e engenheiros econômicos.

\section{Digressão e Vitríolo}

A teoria e a evidência empírica do crescimento econômico de longo prazo estão além do escopo do ensaio, mas vale a pena notar que esses tópicos ocuparam grande parte da atenção dos macroeconomistas durante a década de 1990. Este trabalho desviou a atenção das flutuações de curto prazo, que dominaram o campo da macroeconomia desde o seu nascimento, meio século antes.

Há várias razões para o surgimento do crescimento como uma grande área de pesquisa. Primeiro, uma série de artigos influentes de Paul Romer (1986) e outros ofereceram um novo conjunto de ideias e ferramentas para analisar o que é certamente um dos tópicos mais instigantes da economia — a enorme lacuna entre nações ricas e pobres. Segundo, novos dados comparativos entre países tornaram-se disponíveis e permitiram o exame sistemático da validade de teorias alternativas (Summers e Heston, 1991). Terceiro, a economia dos EUA na década de 1990 estava vivenciando sua expansão mais longa na história. Assim como os primeiros keynesianos foram atraídos para o campo por causa de sua relevância imediata para a saúde da nação, a economia da década de 1990 sugeriu àquela geração de estudantes que o ciclo econômico não era mais de grande importância prática.

Há também uma quarta razão, mais problemática, pela qual os macroeconomistas iniciantes da década de 1990 foram levados a estudar o crescimento de longo prazo em vez das flutuações de curto prazo: a tensão entre as visões de mundo nova clássica e nova keynesiana. Enquanto Lucas, o principal economista novo clássico, proclamava que "as pessoas não levam mais a sério a teorização keynesiana", keynesianos proeminentes eram igualmente desdenhosos com seus colegas novos clássicos. Em seu Discurso Presidencial para a American Economic Association, Solow (1980) chamou de "tolamente restritivo" o fato de os novos economistas clássicos descartarem por suposição a existência de rigidez de salários e preços e a possibilidade de os mercados não se equilibrarem. Ele disse: "Lembro-me de ler uma vez que ainda não se entende como a girafa consegue bombear um suprimento de sangue adequado até sua cabeça; mas é difícil imaginar que alguém, portanto, concluiria que as girafas não têm pescoços longos".

Em uma entrevista com Arjo Klamer (1984) alguns anos depois, Lucas comentou: "Não acho que Solow, em particular, tenha tentado lidar com qualquer uma dessas questões exceto fazendo piadas". Em sua própria entrevista no mesmo volume, Solow explicou sua indisposição em se envolver com os novos economistas clássicos: "Suponha que alguém se sente onde você está sentado agora e anuncie para mim que ele é Napoleão Bonaparte. A última coisa que eu quero fazer com ele é me envolver em uma discussão técnica sobre táticas de cavalaria na Batalha de Austerlitz. Se eu fizer isso, estarei sendo tacitamente atraído para o jogo de que ele é Napoleão Bonaparte".

Até certo ponto, essa disputa reflete as diferentes perspectivas dos protagonistas sobre o objetivo do campo. Lucas parece estar reclamando que Solow não aprecia o maior rigor analítico que a macroeconomia nova clássica pode oferecer. Solow parece estar reclamando que Lucas não aprecia a patente falta de realidade de suas suposições de equilíbrio de mercado. Cada um deles tem um ponto. Do ponto de vista da ciência, o maior rigor que os novos clássicos ofereciam tem muito apelo. Mas, do ponto de vista da engenharia, o custo desse rigor adicional parece pesado demais para suportar.

Eu me detenho na natureza desse debate não apenas porque reflete a tensão subjacente entre cientistas e engenheiros, mas também porque ajuda a explicar as escolhas feitas pela próxima geração de economistas. Tal vitríolo entre gigantes intelectuais atrai a atenção, da mesma forma que os frequentadores de um bar se reúnem em torno de uma briga, incentivando os participantes. Mas não foi saudável para o campo da macroeconomia. Sem surpresa, muitos jovens economistas optaram por evitar tomar partido nessa disputa, voltando sua atenção para longe das flutuações econômicas e em direção a outros tópicos.

\section{Uma Nova Síntese, ou uma Trégua?}

Um velho ditado diz que a ciência progride funeral por funeral. Hoje, com os benefícios de uma maior expectativa de vida, seria mais preciso (se menos vívido) dizer que a ciência progride aposentadoria por aposentadoria. Na macroeconomia, à medida que a geração mais velha de protagonistas se aposentou ou se aproximou da aposentadoria, foi substituída por uma geração mais jovem de macroeconomistas que adotaram uma cultura de maior civilidade. Ao mesmo tempo, um novo consenso surgiu sobre a melhor maneira de entender as flutuações econômicas. Marvin Goodfriend e Robert King (1997) apelidaram essa visão de consenso de "a nova síntese neoclássica". Este modelo de síntese tem sido amplamente aplicado na pesquisa sobre política monetária (Clarida, Gali e Gertler, 1999; McCallum e Nelson, 1999). O tratamento mais extenso desta nova síntese é o tratado monumental (em ambos os sentidos da palavra) de Michael Woodford (2003).

Como a síntese neoclássica-keynesiana de uma geração anterior, a nova síntese tenta fundir os pontos fortes das abordagens concorrentes que a precederam. Dos modelos novos clássicos, ela toma as ferramentas da teoria do equilíbrio geral estocástico dinâmico. Preferências, restrições e otimização são o ponto de partida, e a análise é construída a partir desses microfundamentos. Dos modelos novos keynesianos, ela toma as rigidezes nominais e as usa para explicar por que a política monetária tem efeitos reais no curto prazo. A abordagem mais comum é assumir empresas monopolisticamente competitivas que mudam os preços apenas intermitentemente, resultando em uma dinâmica de preços às vezes chamada de nova curva de Phillips keynesiana. O cerne da síntese é a visão de que a economia é um sistema de equilíbrio geral dinâmico que se desvia de um ótimo de Pareto por causa de preços rígidos (e talvez uma variedade de outras imperfeições de mercado).

É tentador descrever o surgimento desse consenso como um grande progresso. De certa forma, é. Mas também há uma maneira menos otimista de ver o estado atual das coisas. Talvez o que tenha ocorrido não tenha sido tanto uma síntese, mas uma trégua entre combatentes intelectuais, seguida de um recuo para salvar as aparências de ambos os lados. Tanto novos clássicos quanto novos keynesianos podem olhar para essa nova síntese e reivindicar um grau de vitória, ignorando a derrota mais profunda que jaz sob a superfície.

O coração dessa nova síntese — um sistema de equilíbrio geral dinâmico com rigidezes nominais — é precisamente o que se encontra nos primeiros modelos keynesianos. Hicks propôs o modelo IS-LM, por exemplo, em uma tentativa de colocar as ideias de Keynes em um cenário de equilíbrio geral. (Lembre-se de que Hicks ganhou o Prêmio Nobel de 1972 juntamente com Kenneth Arrow por contribuições à teoria do equilíbrio geral.) Klein, Modigliani e os outros construtores de modelos estavam tentando levar esse sistema de equilíbrio geral aos dados para conceber políticas melhores. Em grande medida, a nova síntese retoma a agenda de pesquisa que a profissão abandonou, a mando dos novos clássicos, na década de 1970.

Com o benefício da retrospectiva, fica claro que os economistas novos clássicos prometeram mais do que podiam entregar. Seu objetivo declarado era descartar a teorização keynesiana e substituí-la por modelos de equilíbrio de mercado que pudessem ser convincentemente levados aos dados e então usados para análise de políticas. Por esse padrão, o movimento falhou. Em vez disso, ajudaram a desenvolver ferramentas analíticas que agora estão sendo usadas para desenvolver outra geração de modelos que assumem preços rígidos e que, de muitas maneiras, se assemelham aos modelos contra os quais os novos clássicos estavam em campanha.

Os novos keynesianos podem reivindicar um grau de vindicação aqui. A nova síntese descarta a suposição de equilíbrio de mercado que Solow chamou de "tolamente restritiva" e que a pesquisa nova keynesiana sobre preços rígidos visava minar. No entanto, os novos keynesianos podem ser criticados por terem caído na isca dos novos clássicos e, como resultado, buscado um programa de pesquisa que acabou sendo abstrato demais e insuficientemente prático. Paul Krugman (2000) oferece esta avaliação do programa de pesquisa novo keynesiano: "Pode-se agora explicar como a rigidez de preços \textit{poderia} acontecer. Mas previsões úteis sobre quando ela acontece e quando não acontece, ou modelos que partam dos custos de menu para uma curva de Phillips realista, simplesmente parecem não estar surgindo". Mesmo como proponente desta linha de trabalho, tenho que admitir que há alguma verdade nessa avaliação.

\section{A Visão do Banco Central}

Se Deus colocou macroeconomistas na Terra para resolver problemas práticos, então São Pedro acabará por nos julgar por nossas contribuições à engenharia econômica. Então vamos perguntar: os desenvolvimentos na teoria do ciclo econômico nas últimas décadas melhoraram a formulação da política econômica? Ou, para estabelecer um objetivo mais modesto, os avanços na ciência macroeconômica alteraram a forma como a política econômica é analisada e discutida entre os economistas profissionais envolvidos no processo político?

Um lugar para encontrar evidências para responder a essas perguntas é o encantador livro de memórias de Laurence Meyer, \textit{A Term at the Fed}. Em 1996, Meyer deixou seu emprego como professor de economia na Washington University e como proeminente consultor econômico para servir por seis anos como governador do Federal Reserve. Seu livro oferece uma janela sobre como os economistas nos escalões mais altos da formulação de política monetária veem seus trabalhos e as abordagens que adotam para analisar a economia.

O livro deixa o leitor com uma impressão clara: os desenvolvimentos recentes na teoria do ciclo econômico, promulgados tanto por novos clássicos quanto por novos keynesianos, tiveram um impacto próximo de zero na formulação prática de políticas. A análise de Meyer sobre as flutuações econômicas e a política monetária é inteligente e nuançada, mas não mostra traços da teoria macroeconômica moderna. Ela pareceria quase completamente familiar para alguém que foi formado na síntese neoclássica-keynesiana que prevalecia por volta de 1970 e ignorou a literatura acadêmica desde então. A visão de mundo de Meyer seria fácil de descartar como ultrapassada se fosse idiossincrática, mas não é. Ela é típica dos economistas que ocuparam cargos de destaque nos bancos centrais do mundo.

É moda entre os acadêmicos acreditar que os bancos centrais foram fortemente influenciados pela literatura sobre regras versus discrição, particularmente o trabalho sobre inconsistência temporal que começou com Kydland e Prescott (1977). Duas mudanças institucionais são frequentemente ligadas a essas contribuições acadêmicas: a maior independência dos bancos centrais em países como a Nova Zelândia e a adoção de metas de inflação como regime de política em muitos bancos centrais ao redor do mundo. Essas mudanças institucionais, por sua vez, são ligadas a melhorias na política monetária. De acordo com essa linha de argumentação, Kydland e Prescott devem ser agradecidos pela inflação baixa e estável que muitos países desfrutaram nas últimas duas décadas.

Essa visão autoelogiosa depara-se com dois problemas. Primeiro, as mudanças institucionais que observamos estão, no máximo, vagamente ligadas às questões levantadas na literatura teórica. Um banco central independente não é o mesmo que um banco central vinculado a regras. O Federal Reserve dos EUA há muito tempo tem um alto grau de independência sem nunca se comprometer com uma regra de política. Mesmo as metas de inflação estão mais próximas de uma declaração de intenções e uma forma de comunicação com o público do que de um compromisso com uma regra de política. Bernanke (2003) chamou isso de "discrição restringida".

O segundo problema, mais significativo, é que essas mudanças institucionais não estão necessariamente ligadas às melhorias que testemunhamos na política monetária. Ball e Sheridan (2005) analisam uma grande amostra de países e mostram que a adoção de metas de inflação não ajuda a explicar o movimento recente em direção a uma inflação baixa e estável. A política monetária melhorou tanto nos países que adotaram metas de inflação quanto nos que não as adotaram. Essa melhoria mundial nos resultados da inflação pode ser porque a economia mundial não teve que lidar com choques de oferta tão adversos quanto os experimentados na década de 1970 ou porque os banqueiros centrais aprenderam com a experiência da década de 1970 que a inflação alta deve ser assiduamente evitada. Mas as evidências mostram que as metas de inflação não são um pré-requisito para uma boa política monetária.

O Federal Reserve sob Alan Greenspan é um caso exemplar. De acordo com Blinder e Reis (2005), Greenspan tem o direito de ser chamado de "o maior banqueiro central que já viveu". De fato, segundo a maioria dos relatos, a política monetária funcionou extraordinariamente bem sob sua liderança. No entanto, durante todo o seu tempo à frente do Fed, Greenspan evitou qualquer anúncio de uma regra de política, valorizando a flexibilidade em detrimento do compromisso. Veja como Greenspan (2003) defendeu sua escolha: "Alguns críticos argumentaram que tal abordagem à política é excessivamente indisciplinada — baseada em julgamentos, aparentemente discricionária e difícil de explicar. O Federal Reserve deveria, concluem alguns, tentar ser mais formal em suas operações, vinculando suas ações exclusivamente às prescrições de uma regra de política formal. Que qualquer abordagem nesse sentido levaria a uma melhoria no desempenho econômico, no entanto, é altamente duvidoso... Regras por sua natureza são simples e, quando existem incertezas significativas e mutáveis no ambiente econômico, elas não podem substituir os paradigmas de gerenciamento de risco, que são muito mais adequados à formulação de políticas". No entanto, apesar da aversão de Greenspan a regras de política, a inflação foi baixa e estável durante seu mandato como presidente do Fed. Greenspan prova que os bancos centrais podem produzir resultados desejáveis enquanto exercem poderes discricionários substanciais.

\section{A Visão da Política Fiscal}

Outro lugar para procurar o impacto prático da teoria macroeconômica é a análise da política fiscal. Os cortes de impostos de Bush de 2001 e 2003 oferecem um bom estudo de caso, em parte porque são uma tentativa recente de um grande estímulo fiscal para combater uma recessão e em parte porque, como presidente do Conselho de Consultores Econômicos por dois anos, estou familiarizado com grande parte da análise econômica que lançou as bases para essa política. Certamente, houve muitos motivos para o desenho da política tributária de Bush. A expansão do crédito tributário infantil, por exemplo, estava enraizada tanto na política e na filosofia social quanto na economia. Mas os economistas do Conselho de Consultores Econômicos e do Tesouro tiveram uma contribuição substancial no desenvolvimento da política, por isso é esclarecedor considerar as ferramentas que eles trouxeram para o trabalho.

A análise econômica do plano tributário de Bush foi feita com um olho no crescimento de longo prazo e o outro no ciclo econômico de curto prazo. A perspectiva de longo prazo seria familiar para estudantes de finanças públicas. Mais significativamente, em 2003, Bush propôs eliminar a dupla tributação da renda do capital corporativo. O projeto final aprovado pelo Congresso não alcançou totalmente esse objetivo, mas o corte substancial nas alíquotas de impostos sobre dividendos moveu-se na direção de uma maior neutralidade tributária, reduzindo o viés para lucros retidos sobre dividendos, o viés para o financiamento por dívida sobre o capital próprio e o viés para o capital não corporativo sobre o corporativo. Também moveu o código tributário mais na direção de tributar o consumo em vez da renda. Este último objetivo é consistente com uma literatura bem estabelecida em finanças públicas (por exemplo, Diamond e Mirrlees, 1971; Atkinson e Stiglitz, 1976; Feldstein, 1978; Chamley, 1986) e não é particularmente novo como uma questão de teoria econômica. Há três décadas, Atkinson e Stiglitz notaram que, mesmo então, havia uma "presunção convencional em favor da tributação do consumo em vez da renda".

Mais relevante para este ensaio, no entanto, é a análise de curto prazo da política tributária. Quando o presidente George W. Bush assumiu o cargo em 2001, a economia estava entrando em uma recessão após o estouro da bolha do mercado de ações do final da década de 1990. Um objetivo dos cortes de impostos era estimular a recuperação econômica e o emprego. Quando Bush assinou o \textit{Jobs and Growth Tax Relief Reconciliation Act} de 2003, explicou a política da seguinte forma: "Quando as pessoas têm mais dinheiro, elas podem gastá-lo em bens e serviços. E em nossa sociedade, quando elas demandam um bem ou serviço adicional, alguém produzirá o bem ou serviço. E quando alguém produz esse bem ou serviço, significa que alguém tem mais probabilidade de conseguir um emprego". Essa lógica é quintessencialmente keynesiana.

O Conselho de Consultores Econômicos foi solicitado a quantificar como o alívio tributário afetaria o emprego. Respondemos a essa pergunta usando um modelo macroeconométrico convencional. O modelo específico que usamos enquanto eu estava lá era o mantido pela Macroeconomic Advisers, a consultoria criada e dirigida por Laurence Meyer antes de ele ser um governador do Federal Reserve. Este modelo estava sendo usado pela equipe do Conselho de Consultores Econômicos muito antes de eu chegar como presidente e, de fato, tinha sido usado por quase duas décadas sob as administrações republicana e democrata. A escolha deste modelo em particular não é crucial, no entanto, pois o modelo da Macroeconomic Advisers é semelhante a outros grandes modelos macroeconométricos, como o modelo FRB/US mantido pelo Federal Reserve. Do ponto de vista da história intelectual, esses modelos são descendentes diretos dos esforços iniciais de modelagem de Klein, Modigliani e Eckstein. A pesquisa dos novos clássicos e dos novos keynesianos teve influência mínima na construção desses modelos.

O mundo real da formulação de políticas macroeconômicas pode ser desanimador para aqueles de nós que passaram a maior parte de nossas carreiras na academia. A triste verdade é que a pesquisa macroeconômica das últimas três décadas teve apenas um impacto menor na análise prática da política monetária ou fiscal. A explicação não é que os economistas na arena política ignoram os desenvolvimentos recentes. Pelo contrário: a equipe do Federal Reserve inclui alguns dos melhores jovens doutores, e o Conselho de Consultores Econômicos, tanto sob as administrações democrata quanto republicana, atrai talentos das melhores universidades de pesquisa do país. O fato de que a pesquisa macroeconômica moderna não é amplamente utilizada na formulação prática de políticas é evidência \textit{prima facie} de que ela é de pouca utilidade para esse propósito. A pesquisa pode ter sido bem-sucedida como uma questão de ciência, mas não contribuiu significativamente para a engenharia macroeconômica.

\section{Dentro da Sala de Aula}

Além dos corredores do poder nas capitais mundiais, há outro lugar onde a profissão econômica tenta vender seus produtos a um público mais amplo — a sala de aula da graduação. Aqueles de nós que regularmente ensinam alunos de graduação veem nosso trabalho como a produção de cidadãos bem informados sobre os princípios de uma boa política. Nossa escolha de material é guiada pelo que vemos como importante para a próxima geração de eleitores entender.

Como os formuladores de políticas, os estudantes de graduação normalmente têm pouco interesse na teoria pela teoria em si. Em vez disso, eles estão interessados em entender como o mundo real funciona e como a política pública pode melhorar o desempenho econômico. Exceto pelo raro estudante que está considerando a pós-graduação e uma carreira como economista acadêmico, o estudante de graduação tem a perspectiva de um engenheiro, mais do que a de um cientista. É, portanto, útil observar o que escolhemos ensinar aos alunos de graduação. E não há lugar melhor para ver o que ensinamos do que no conteúdo dos livros didáticos de graduação mais amplamente utilizados.

Considere, por exemplo, os livros usados para ensinar macroeconomia de nível intermediário. Uma geração atrás, os três principais textos para este curso eram os de Robert Gordon; Robert Hall e John Taylor; e Rudiger Dornbusch e Stanley Fischer. Hoje, os três mais vendidos são os escritos por Olivier Blanchard; Andrew Abel e Ben Bernanke; e eu mesmo. O fio condutor é que cada um desses seis livros foi escrito por pelo menos um economista com treinamento de pós-graduação no MIT, uma escola de engenharia proeminente onde a tradição macroeconômica dominante era a de Samuelson e Solow. Em todos esses livros, a teoria básica ensinada aos alunos de graduação é alguma versão da oferta agregada e demanda agregada, e a teoria básica da demanda agregada é o modelo IS-LM. A mesma lição pode ser colhida ao analisar os livros didáticos mais amplamente usados para a economia do primeiro ano: as flutuações econômicas de curto prazo são melhor compreendidas usando alguma versão da síntese neoclássica-keynesiana.

Não quero sugerir que a pedagogia tenha permanecido estagnada à medida que o campo evoluiu. Os livros didáticos de hoje dão maior ênfase à teoria monetária clássica, aos modelos de crescimento de longo prazo e ao papel das expectativas do que os de 30 anos atrás. Há menos confiança sobre o que a política pode realizar e mais ênfase nas regras de política sobre ações monetárias e fiscais discricionárias (apesar da falta de evidências sobre a importância prática das regras de política). Mas o arcabouço básico que os estudantes modernos aprendem para dar sentido ao ciclo econômico é um que seria familiar a uma geração anterior de keynesianos.

A exceção que prova a regra é o texto intermediário escrito por Robert Barro, publicado pela primeira vez em 1984. O livro de Barro forneceu uma introdução clara e acessível à abordagem nova clássica da macroeconomia voltada para estudantes de graduação. Modelos keynesianos foram incluídos, mas foram cobertos tardiamente no livro, brevemente e com pouca ênfase. Quando o livro saiu, recebeu atenção e aclamação substanciais. No entanto, embora muitos macroeconomistas tenham lido o livro de Barro e ficado impressionados com ele, muitos menos o escolheram para seus alunos. A revolução nova clássica na pedagogia que Barro esperava inspirar nunca decolou, e o texto de Barro não ofereceu competição significativa aos livros didáticos dominantes da época.

Esta falta de revolução na pedagogia macroeconômica contrasta fortemente com o que ocorreu há meio século. Quando o livro introdutório de Samuelson foi publicado pela primeira vez em 1948 com o objetivo de introduzir os estudantes de graduação à revolução keynesiana, os professores do mundo abraçaram rápida e cordialmente a nova abordagem. Em contraste, as ideias dos novos clássicos e dos novos keynesianos não mudaram fundamentalmente a forma como a macroeconomia da graduação é ensinada.

\section{Nenhum Dentista à Vista}

John Maynard Keynes (1931) famosamente opinou: "Se os economistas conseguissem ser pensados como pessoas humildes e competentes em um nível comparável ao dos dentistas, isso seria esplêndido". Ele estava expressando a esperança de que a ciência da macroeconomia evoluísse para um tipo de engenharia útil e rotineiro. Nesse futuro utópico, evitar uma recessão seria tão simples quanto tratar uma cárie.

Os principais desenvolvimentos na macroeconomia acadêmica das últimas décadas têm pouca semelhança com a odontologia. A pesquisa nova clássica e nova keynesiana teve pouco impacto sobre os macroeconomistas práticos encarregados da tarefa confusa de conduzir a política monetária e fiscal real. Também teve pouco impacto sobre o que os professores dizem aos futuros eleitores sobre a política macroeconômica quando entram na sala de aula da graduação. Do ponto de vista da engenharia macroeconômica, o trabalho das últimas décadas parece um infeliz erro de percurso.

No entanto, da perspectiva mais abstrata da ciência macroeconômica, este trabalho pode ser visto de forma mais positiva. Os economistas novos clássicos foram bem-sucedidos em mostrar as limitações dos grandes modelos macroeconométricos keynesianos e as prescrições de políticas baseadas nesses modelos. Eles chamaram a atenção para a importância das expectativas e o caso para regras de política. Os economistas novos keynesianos forneceram modelos melhores para explicar por que os salários e preços não equilibram os mercados e, mais genericamente, que tipos de imperfeições de mercado são necessários para dar sentido às flutuações econômicas de curto prazo. A tensão entre essas duas visões, embora nem sempre civilizada, pode ter sido produtiva, pois a competição é tão importante para o avanço intelectual quanto para os resultados do mercado.

Os \textit{insights} resultantes estão sendo incorporados na nova síntese que agora está se desenvolvendo e que, eventualmente, se tornará a base para a próxima geração de modelos macroeconométricos. Para aqueles de nós interessados na macroeconomia como ciência e engenharia, podemos ver o surgimento recente de uma nova síntese como um sinal de esperança de que mais progresso possa ser feito em ambas as frentes. À medida que olhamos para o futuro, "humilde" e "competente" permanecem ideais aos quais os macroeconomistas podem aspirar.

\section*{Referências}

\begin{itemize}
    \item Akerlof, George A. e Janet L. Yellen. 1985. "A Near-Rational Model of the Business Cycle with Wage and Price Inertia." \textit{Quarterly Journal of Economics}. 100:5 (Suplemento), pp. 823–38.
    \item Atkinson, A. B. e J. E. Stigltiz. 1976. "The Design of Tax Structure: Direct versus Indirect Taxation," \textit{Journal of Public Economics}. 6:1–2, pp. 55–75.
    \item Ball, Laurence e David Romer. 1990. "Real Rigidities and the Non-Neutrality of Money." \textit{Review of Economic Studies}. 57:2, pp. 183–203.
    \item Ball, Laurence e Niamh Sheridan. 2005. "Does Inflation Targeting Matter?" em \textit{The Inflation-Targeting Debate}. Ben S. Bernanke e Michael Woodford, eds. Chicago: University of Chicago Press, pp. 249–82.
    \item Barro, Robert. 1977. "Unanticipated Money Growth and Unemployment in the United States." \textit{American Economic Review}. Março, 67:2, pp. 101–15.
    \item Barro, Robert J. e Herschel Grossman. 1971. "A General Disequilibrium Model of Income and Employment." \textit{American Economic Review}. Março, 61:1, pp. 82–93.
    \item Bernanke, Ben S. 2003. "Constrained Discretion and Monetary Policy." Discurso perante os Money Marketeers da New York University, Nova York, Nova York. 3 de fevereiro.
    \item Blanchard, Olivier Jean, e Nobuhiro Kiyotaki. 1987. "Monopolistic Competition and the Effects of Aggregate Demand." \textit{American Economic Review}. Setembro, 77:4, pp. 647–66.
    \item Blinder, Alan S. e Ricardo Reis. 2005. "Understanding the Greenspan Standard," Artigo apresentado na conferência do Federal Reserve Bank of Kansas City em Jackson Hole. (Princeton University.)
    \item Breit, William e Barry T. Hirsch. 2004. \textit{Lives of the Laureates}, 4ª edição. Cambridge, MA: MIT Press.
    \item Chamley, Christophe. 1986. "Optimal Taxation of Capital Income in General Equilibrium with Infinite Lives." \textit{Econometrica}. Maio, 54:3, pp. 607–22.
    \item Clarida, Richard, Jordi Gali e Mark Gertler. 1999. "The Science of Monetary Policy: A New Keynesian Perspective." \textit{Journal of Economic Literature}. Dezembro, 37:4, pp. 1661–1707.
    \item Diamond, Peter A. e James A. Mirrlees. 1971. "Optimal Taxation and Public Production II: Tax Rules. \textit{American Economic Review}. Junho, 61:3, pp. 261–78.
    \item Feldstein, Martin. 1978. "The Welfare Cost of Capital Income Taxation," \textit{Journal of Political Economy}. Abril, 86:2, pt. 2, pp. S29–S51.
    \item Fischer, Stanley. 1977. "Long-term Contracts, Rational Expectations, and the Optimal Money Supply Rule." \textit{Journal of Political Economy}. Fevereiro, 85:1, pp. 191–205.
    \item Friedman, Milton. 1957. \textit{A Theory of the Consumption Function}. Princeton: Princeton University Press.
    \item Friedman, Milton. 1968. "The Role of Monetary Policy." \textit{American Economic Review}. Março, 58:1, pp. 1–17.
    \item Friedman, Milton e Anna Jacobson Schwartz. 1963. \textit{A Monetary History of the United States, 1867–1960}. Princeton: Princeton University Press.
    \item Garfield, Eugene. 1990. "Who Will Win the Nobel Prize in Economics? Here’s a Forecast Based on Citation Indicators," em \textit{Essays of an Information Scientist: Journalology, KeyWords Plus, and Other Essays}, Vol. 13. Philadelphia, PA: ISS Press, p. 83–87.
    \item Goodfriend, Marvin e Robert King. 1997. "The New Neoclassical Synthesis and the Role of Monetary Policy," em \textit{NBER Macroeconomics Annual}. Ben S. Bernanke e Julio J. Rotemberg, eds. Cambridge, MA: MIT Press. pp. 231–83.
    \item Greenspan, Alan. 2003. "Monetary Policy under Uncertainty." Discurso em simpósio patrocinado pelo Federal Reserve Bank of Kansas City, Jackson Hole, Wyoming, 29 de agosto.
    \item Hicks, John R. 1937. "Mr. Keynes and the ‘Classics’." \textit{Econometrica}. Abril, 5:2, pp. 147–59.
    \item Hume, David. 1752. "Of Money," em \textit{Essays}. Londres: George Routledge and Sons.
    \item Keynes, John Maynard. 1931. \textit{Essays in Persuasion}. Nova York: Norton.
    \item Klamer, Arjo. 1984. \textit{Conversations with Economists}. Totowa, NJ: Rowman and Allanheld.
    \item Klein, Lawrence R. 1946. "Macroeconomics and the Theory of Rational Behavior." \textit{Econometrica}. Abril, 14:2, pp. 93–108.
    \item Krugman, Paul. 2000. "How Complicated Does the Model Have to Be?" \textit{Oxford Review of Economic Policy}. 16:4, pp. 33–42.
    \item Kydland, Finn e Edward C. Prescott. 1977. "Rules Rather Than Discretion: The Inconsistency of Optimal Plans." \textit{Journal of Political Economy}. Junho, 85:3, pp. 473–92.
    \item Kydland, Finn e Edward C. Prescott. 1982. "Time to Build and Aggregate Fluctuations." \textit{Econometrica}. Novembro, 50:6, pp. 1345–71.
    \item Long, John B. e Charles Plosser. 1983. "Real Business Cycles." \textit{Journal of Political Economy}. Fevereiro, 91:1, pp. 39–69.
    \item Lucas, Robert E., Jr. 1973. "Some International Evidence on Output–Inflation Tradeoffs." \textit{American Economic Review}. Junho, 63:3, pp. 326–34.
    \item Lucas, Robert E. Jr. 1976. "Econometric Policy Evaluation: A Critique." \textit{Carnegie-Rochester Conference Series on Public Policy}. 1, pp. 19–46.
    \item Lucas, Robert E., Jr. 1980. "The Death of Keynesian Economics." \textit{Issues and Ideas}. (University of Chicago, Chicago, IL). Inverno, pp. 18–19.
    \item Lucas, Robert E. Jr. e Thomas J. Sargent. 1979. "After Keynesian Macroeconomics." \textit{Federal Reserve Bank of Minneapolis Quarterly Review}. Primavera, 3:2, pp. 1–16.
    \item Malinvaud, Edmund. 1977. \textit{The Theory of Unemployment Reconsidered}. Oxford, UK: Blackwell.
    \item Mankiw, N. Gregory. 1985. "Small Menu Costs and Large Business Cycles: A Macroeconomic Model of Monopoly." \textit{Quarterly Journal of Economics}. Maio, 100:2, pp. 529–37.
    \item McCallum, Bennett T. e Edward Nelson. 1999. "An Optimizing IS-LM Specification for Monetary Policy and Business Cycle Analysis." \textit{Journal of Money, Credit, and Banking}. Agosto, 31:3, pt. 1, pp. 296–316.
    \item Meyer, Laurence H. 2004. \textit{A Term at the Fed: An Insider’s View}. Nova York: Harper-Collins.
    \item Modigliani, Franco. 1944. "Liquidity Preference and the Theory of Interest and Money." \textit{Econometrica}. Janeiro, 12:1, pp. 45–88.
    \item Phelps, Edmund. 1968. "Money Wage Dynamics and Labor Market Equilibrium." \textit{Journal of Political Economy}. Julho/Agosto, 76:4, pt. 2, pp. 678–711.
    \item Pigou, Arthur. 1927. \textit{Industrial Fluctuations}. Londres: Macmillan and Company.
    \item Quandt, Richard E. 1976. "Some Quantitative Aspects of the Economics Journal Literature." \textit{Journal of Political Economy}. Agosto, 84:4, pt. 1, pp. 741–55.
    \item Romer, Paul. 1986. "Increasing Returns and Long Run Growth." \textit{Journal of Political Economy}. Outubro, 94:5, pp. 1002–37.
    \item Samuelson, Paul A. 1948. \textit{Economics: An Introductory Analysis}. Nova York: McGraw-Hill.
    \item Samuelson, Paul A. 1988. "Keynesian Economics and Harvard: In the Beginning." \textit{Challenge: The Magazine of Economic Affairs}. Julho/Agosto, 31, pp. 32–34.
    \item Samuelson, Paul A. e Robert M. Solow. 1960. "Analytical Aspects of Anti-Inflation Policy." \textit{American Economic Review}. Maio, 50:2, pp. 177–194.
    \item Sargent, Thomas J. e Neil Wallace. 1975. "Rational Expectations, the Optimal Monetary Instrument, and the Optimal Money Supply Rule." \textit{Journal of Political Economy}. Abril, 83:2, pp. 241–54.
    \item Solow, Robert M. 1980. "On Theories of Unemployment." \textit{American Economic Review}. Março, 70:1, pp. 1–11.
    \item Summers, Robert e Alan Heston. 1991. "The Penn World Table (Mark 5): An Expanded Set of International Comparisons, 1950–1988." \textit{Quarterly Journal of Economics}. Maio, 106:2, pp. 327–68.
    \item Taylor, John B. 1980. "Aggregate Dynamics and Staggered Contracts." \textit{Journal of Political Economy}. Fevereiro, 88:1, pp. 1–23.
    \item Woodford, Michael. 2003. \textit{Interest and Prices}. Princeton: Princeton University Press.
\end{itemize}

\end{document}